\band{orangehead}{Lecture 3 · Specifications / SDD / Tests}
\topic{SRS → SDS → STS}
\k{SRS}: ต้องการอะไร; \k{SDS}: ออกแบบอย่างไร; \k{STS}: ตรวจอะไร ด้วยข้อมูลใด คาดผลอย่างไร และเก็บหลักฐานใด
\b{SRS → feature inventory → req-feature map → system SDS → sprint vertical slice → feature contract → implement → test/review → integrate → main}
\k{System SDS}: โครงสร้าง/stack/data, auth, มาตรฐาน API/UI/errors และ NFR ร่วม
\k{Feature SDS}: ขอบเขต/req, workflow/state/BR, data/API/UI, สิทธิ์/validation/errors, AC/tests
\k{Contract 7}: \b{inputs → validations → behavior → persistence → side effects → errors → outputs}
\g{Invoice}: รับ customer/date/lines → ตรวจ customer, qty > 0, price ≥ 0 → คำนวณ → บันทึกหัว/รายการใน transaction เดียว; ผิด → ไม่บันทึก \extra receipt เป็น side effect ที่ต้องตกลง
\topic{Tests / Red-Green-Refactor}
\b{RED}: test ล้มเหลวด้วยเหตุผลที่ตั้งใจ → \b{GREEN}: เขียนโค้ดเท่าที่จำเป็นให้ผ่าน → \b{REFACTOR}: ปรับโครงสร้างโดย tests ยังผ่าน
\k{AAA}: Arrange เตรียมข้อมูล → Act เรียกใช้ → Assert ตรวจผล; ครอบคลุมปกติ/ผิดเงื่อนไข/ขอบเขต/error/สิทธิ์; \r{ต้องจับ defect ได้ ไม่เขียนตาม implementation อย่างเดียว}
\k{STS row}: test ID + Req/AC → precondition/data → steps/input → expected UI/API/DB → file + final result
\extra \g{ค่าขอบเขต (boundary)}: คะแนน 0-100 → ทดสอบ −1/0/100/101; 0 และ 100 ผ่าน ส่วน −1 และ 101 ถูกปฏิเสธ
{\smalltable\begin{tabularx}{\linewidth}{|p{11mm}|Y|}\hline
\rowcolor{lightrow}ระดับ/ชนิด & ตรวจอะไร\\\hline
Unit & function/class/component; แยก deps ด้วย mock\\\hline
Integration & modules/API/DB/interfaces ทำงานร่วมกันจริง\\\hline
System & ระบบรวมตรง SRS ทั้ง FR/NFR\\\hline
E2E & เส้นทางผู้ใช้ผ่าน UI → API → DB/external\\\hline
UAT & ผู้ใช้/ธุรกิจยอมรับตาม need/AC และ sign-off\\\hline
Regression & ทดสอบของเดิมหลังแก้ไข; ใช้ได้ทุกระดับ\\\hline
Smoke & เปิดระบบและใช้งานหลักได้หลัง setup/release\\\hline
Special & load/performance, security, compatibility, recovery\\\hline
\end{tabularx}}
\k{Verification}: ทำตรง spec ไหม; \k{Validation}: ตรงความต้องการผู้ใช้ไหม; \r{coverage สูง ≠ tests จับข้อผิดพลาดได้ดี}
\k{Sync}: รอผลตอบ; \k{Async}: ส่งงานแล้วรับผลภายหลัง; ระบุการลองซ้ำ (retry) และ idempotency: ส่งซ้ำแล้วไม่เกิดผลซ้ำเกินที่ตกลง
\k{POS recovery/outbox}: เก็บ payment attempt; ผลไม่ชัดตรวจ PSP ก่อนเก็บเงินซ้ำ; บันทึก sale + pending inventory แบบ atomic; ส่ง async รอ ack/retry ใช้ transaction ID กัน stock ซ้ำ
\k{BDD}: เริ่มจากพฤติกรรมด้วย Gherkin (outside-in); \k{TDD}: ให้ tests นำการออกแบบ (inside-out); ใช้ร่วมกันได้
\chapterband{orangehead}{Lecture 4 · UML / Feature Design / UI}
\k{Use case}: เป้าหมายของ actor ภายใน system boundary; วงรี = goal; actor association \r{เส้นทึบ ไม่มีหัวลูกศร}
\b{include}: base → ขั้นตอนบังคับ (เส้นประ); \b{extend}: กรณีเสริม → base (เส้นประ + condition); \b{generalization}: child → parent (สามเหลี่ยมโปร่ง)
\begin{center}\begin{tikzpicture}[x=1mm,y=1mm,tinytext]
\node[draw,ellipse,inner sep=1.3pt](a)at(10,0){Checkout};
\node[draw,ellipse,inner sep=1.3pt](b)at(36,0){Pay};
\draw[dependency](a)--node[above,font=\fontsize{6.5}{6.8}\selectfont]{\textless\textless include\textgreater\textgreater}(b);
\node[draw,ellipse,inner sep=1.3pt](c)at(10,-6){Discount};
\draw[dependency](c)--node[right]{\textless\textless extend\textgreater\textgreater}(a);
\end{tikzpicture}\end{center}
\k{Activity}: workflow ครบทางเลือก; ● เริ่ม, action มุมมน, ◇ decision/merge + [guard], bar = fork/join ขนาน, ◎ จบ; swimlane แบ่งผู้รับผิดชอบ; ลูกศรข้าม lane ได้
\begin{center}\begin{tikzpicture}[x=1mm,y=1mm,tinytext]
\node[act](a)at(3,-4.5){A};\node[act](b)at(24,-2){B};\node[act](c)at(24,-7){C};\node[act](d)at(45,-4.5){D};
\draw[line width=1.1pt](10,-2)--(10,-7);\draw[line width=1.1pt](38,-2)--(38,-7);
\draw[flow](a)--(10,-4.5);\draw[flow](10,-2)--(b);\draw[flow](10,-7)--(c);\draw[flow](b)--(38,-2);\draw[flow](c)--(38,-7);\draw[flow](38,-4.5)--(d);
\node[anchor=north,font=\fontsize{6.8}{7}\selectfont,inner sep=.3pt]at(24,-8.6){fork: ทำ B/C ขนาน; join: รอครบทั้งคู่};
\end{tikzpicture}\end{center}
\k{Sequence}: scenario หนึ่งใน workflow; participant/lifeline อยู่ด้านบน เวลาลงล่าง; call เส้นทึบ/return เส้นประ; activation คือช่วงทำงาน; \b{alt [condition], opt, loop}; แสดงระบบภายนอกและผลข้างเคียง
\k{State}: สถานะของ \r{object เดียว}; event [guard] / effect; ระบุจุดเริ่ม/จบและ transition ที่อนุญาต; อย่าสับสน action กับ state
\begin{center}\begin{tikzpicture}[x=1mm,y=1mm,tinytext]
\node[state](a)at(6,0){Draft};\node[state](b)at(25,0){Ready};\node[state](c)at(44,0){Submitted};
\draw[flow](a)--node[above]{complete}(b);\draw[flow](b)--node[above]{submit}(c);
\end{tikzpicture}\end{center}
\k{Class (Samples: Final)}: attributes:type; links + multiplicity; inheritance → parent; composition ◆ ที่ whole; aggregation ◇; \r{Samples ห้าม methods}
\k{Scenario}: basic = ปกติ; alternative = ทางเลือกที่ถูกต้อง; exception = ผิดพลาด; Activity ≠ Sequence; Use Case ≠ flowchart
\topic{Generic UI: List / Form / LoV}
\k{List}: title/user/role, search/filter/sort, columns ตาม req, select/open/create, pagination, empty/loading/error; \k{Form}: header + line items + total/status + controls
\k{สัญลักษณ์สอบ}: \r{*} นำหน้าช่องบังคับ; textbox แก้ได้; ขีดเฉียง = read-only/calculated; เลือก FK ด้วย \b{LoV}; รายการย่อยมี add/remove และ validation
\k{ตามสถานะ}: ใน KU Registration, Draft แก้/Save/Submit ได้; Done ล็อก; แสดงผลรอ/สำเร็จ/ผิดพลาดตาม contract; \r{อย่านำชื่อ state นี้ไปใช้ทุกระบบ}
\k{Design system}: สี/ตัวอักษร/ระยะ/components/grid/breakpoints/contrast/keyboard/focus; ตารางยาวอาจเลื่อนแนวนอน
\r{ไม่เติม split payment/cancel/refund หากโจทย์ไม่ให้}; ทุก field/button/transition ต้องโยง req และสิทธิ์ได้
\topic{Case patterns จาก Lecture 4 / review}
{\smalltable\begin{tabularx}{\linewidth}{|p{8mm}|Y|}\hline
\rowcolor{lightrow}Case & สิ่งที่ต้องรักษาเมื่อวาด UML/UI\\\hline
POS & Cashier/Customer; รายการสินค้า/ยอด; cash/card; card ไม่ผ่าน → retry/cash/cancel; สำเร็จ → receipt/stock\\\hline
Hotel & Guest ค้นวัน/เมือง เลือกหลายห้อง/ประเภท จ่าย KBank → ลดห้องว่าง/ยืนยัน/email ถึง Guest และ Hotel; Hotel Admin ดูแลโรงแรมตน\\\hline
Ticket & Customer สร้าง Requested พร้อม action อัตโนมัติ; CCE/Programmer เพิ่ม action ตาม state; action แก้/ลบไม่ได้; latest status/due อัปเดต; ปิดแล้วห้ามเพิ่ม\\\hline
Grades & Instructor เลือก semester/course/sections ดึงรายชื่อ/กรอกเกรด; ครบ → Ready → submit → Submitted; อัปเดต Transcript/email; ใช้ SSO/Registration/Student Info/Email ภายนอก\\\hline
\end{tabularx}}
\k{ตรวจรูป}: ตรง req; guard/ผลข้างเคียงครบ; กรณีผิดกลับแก้ได้; ไม่แต่ง state/สิทธิ์
